\chapter{Hipótese}
\label{cap:hipotese}

A aplicação de análise de sentimento automatizada, utilizando Modelos de Linguagem de
Grande Porte (\gls{LLM}s), como a \gls{API} Gemini, sobre dados textuais gerados por usuários na
plataforma digital Gnomon -- concebida e avaliada sob a metodologia da Design Science --
viabiliza uma identificação mais precisa e profunda da percepção dos usuários sobre os lugares e da
sua construção de identidade territorial. Adicionalmente, postula-se que esta mesma abordagem
analítica possui o potencial de revelar como dimensões de valor socioambiental, que podem se
alinhar a aspectos de Ambientais, Sociais e Governança (\gls{ESG}), são articuladas pelos usuários
em suas experiências e narrativas sobre esses territórios, conforme interagem com a plataforma
Gnomon. A investigação dessas capacidades permitirá o desenvolvimento e a validação de um
método robusto para compreender a experiência do usuário de forma integral, com implicações
diretas para o avanço do conhecimento sobre as intersecções entre tecnologia, identidade de lugar e
a expressão de valores contemporâneos.

Desta forma, a questão central que esta pesquisa busca responder por meio da validação
desta hipótese é:

\begin{citacao}
Como a análise de sentimento aplicada a dados da plataforma Gnomon pode melhor identificar a
percepção das pessoas sobre os lugares e sua identidade de lugar, e, secundariamente, capturar
suas considerações sobre aspectos socioambientais relacionados?
\end{citacao}

\chapter{Objetivo Geral}
\label{cap:objetivo-geral}

O objetivo geral desta pesquisa é investigar, propor e validar uma abordagem metodológica,
fundamentada na Design Science e na análise de sentimento com Modelos de Linguagem de
Grande Porte (LLMs), para identificar e compreender de forma aprimorada a percepção das pessoas
sobre os lugares e a formação de sua identidade de lugar, por meio da interação com a plataforma
digital Gnomon.

\section{Objetivos Específicos}
\label{sec:objetivos-especificos}

Para alcançar o objetivo geral, foram estabelecidos os seguintes objetivos específicos,
adaptados a partir dos focos de investigação de sentimento de pertencimento e percepção de temas
relevantes.

\subsection{Objetivos Específicos Relacionados ao Sentimento de Pertencimento}
\label{subsec:obj-pertencimento}

\begin{enumerate}
    \item \textbf{Identificar Expressões de Sentimento de Pertencimento nos Textos dos Usuários.}
    \begin{alineas}
        \item \textit{Descrição}: Mapear e catalogar palavras-chave, frases e construções linguísticas que indicam um sentimento de pertencimento nos relatos dos usuários.
    \end{alineas}

    \item \textbf{Analisar a Intensidade do Sentimento de Pertencimento.}
    \begin{alineas}
        \item \textit{Descrição}: Utilizar técnicas de análise de sentimentos para medir a intensidade (positiva, neutra, negativa) do sentimento de pertencimento expresso nos textos.
    \end{alineas}

    \item \textbf{Correlacionar o Sentimento de Pertencimento com Fatores Demográficos e Contextuais.}
    \begin{alineas}
        \item \textit{Descrição}: Investigar como variáveis como idade, gênero, origem geográfica influenciam o sentimento de pertencimento.
    \end{alineas}

    \item \textbf{Explorar a Evolução do Sentimento de Pertencimento ao Longo do Tempo.}
    \begin{alineas}
        \item \textit{Descrição}: Estudar como o sentimento de pertencimento dos usuários varia em períodos diferentes, identificando possíveis influências sazonais ou decorrentes de eventos específicos.
    \end{alineas}
\end{enumerate}

\subsection{Objetivos Específicos Relacionados à Percepção dos ODS da ONU}
\label{subsec:obj-ods}

\begin{enumerate}
    \item \textbf{Detectar Menções Diretas e Indiretas aos ODS nos Textos dos Usuários.}
    \begin{alineas}
        \item \textit{Descrição}: Aplicar técnicas de processamento de linguagem natural para identificar referências aos 17 \gls{ODS} nos relatos dos turistas, mesmo quando não mencionados explicitamente.
    \end{alineas}

    \item \textbf{Analisar o Sentimento Associado a Cada ODS Mencionado.}
    \begin{alineas}
        \item \textit{Descrição}: Avaliar a polaridade (positiva, neutra, negativa) das menções aos ODS, identificando quais objetivos são vistos de forma mais favorável ou enfrentam críticas.
    \end{alineas}

    \item \textbf{Avaliar o Nível de Conscientização e Conhecimento dos Indivíduos sobre os ODS.}
    \begin{alineas}
        \item \textit{Descrição}: Investigar o grau de familiaridade dos usuários com os ODS, considerando a profundidade das discussões e a precisão das informações apresentadas nos textos.
    \end{alineas}
\end{enumerate}
